% paper.tex — "Bytes, Not Ratios: What KV Cache Compression Benchmarks Actually Measure"
%
% Portable skeleton: compiles with a stock `article` class so any machine with a LaTeX
% toolchain can build it TODAY. Swap \documentclass for the venue style (neurips_2026.sty /
% iclr2027_conference.sty) before submission — section structure is already venue-agnostic.
%
% FIGURE POLICY (enforced): every figure is included via \datafig, which renders the PDF if and
% only if scripts/ has generated it from results/*.json. Until then it renders a loud placeholder
% naming the script that must run. `make paper` regenerates figures first, so a real build never
% ships a placeholder. No figure is ever hand-drawn or filled with illustrative numbers.

\documentclass[11pt]{article}

\usepackage[utf8]{inputenc}
\usepackage[T1]{fontenc}
\usepackage{microtype}
\usepackage{amsmath,amssymb}
\usepackage{booktabs}
\usepackage{graphicx}
\usepackage{xcolor}
\usepackage{natbib}
\usepackage{hyperref}
\usepackage[margin=1in]{geometry}
\usepackage{tcolorbox}

\hypersetup{colorlinks=true, linkcolor=blue!50!black, citecolor=blue!50!black, urlcolor=blue!50!black}

% \datafig{filename-without-ext}{width}{caption}{label}
% Renders figures/<name>.pdf if present; otherwise a loud red box naming the generating script.
\newcommand{\datafig}[4]{%
  \begin{figure}[t]
    \centering
    \IfFileExists{figures/#1.pdf}{%
      \includegraphics[width=#2\linewidth]{figures/#1.pdf}%
    }{%
      \fbox{\begin{minipage}{#2\linewidth}\centering\vspace{2em}%
        \textcolor{red}{\textbf{[ MISSING DATA FIGURE: figures/#1.pdf ]}}\\[0.5em]
        \small run \texttt{scripts/#1.py} on real \texttt{results/} to generate.\\
        No placeholder numbers are permitted in this slot.\vspace{2em}%
      \end{minipage}}%
    }%
    \caption{#3}
    \label{#4}
  \end{figure}%
}

\title{Bytes, Not Ratios:\\ What KV Cache Compression Benchmarks Actually Measure}
\author{%
  Author name(s) withheld\\
  \texttt{earthespy@gmail.com}%
}
\date{Draft --- \today}

\begin{document}
\maketitle

\begin{abstract}
KV cache compression methods for long-context LLM inference are compared, almost universally, at a
matched \emph{token} budget---a ``compression ratio'' that fixes the fraction of cached tokens
retained---and the result is reported as a memory saving. We show these are not the same quantity,
and that the gap is method-dependent and sometimes total. Auditing the field's standard toolkit, we
find that (i)~some widely-benchmarked methods realize head-wise budgets by \emph{masking} rather than
eviction, so they report large compression ratios while freeing no memory; (ii)~head-level budget
allocation cannot reduce memory under grouped-query attention even when implemented physically,
because keys and values are shared within a query group; and (iii)~peak-memory overheads incurred
during compression are invisible to ratio accounting. We introduce a measurement instrument that
fixes the budget in \emph{realized bytes} and a wrapper that matches any method to a target byte
budget, and re-run the standard comparison. % QUANTITATIVE CLAIMS BELOW ARE FILLED FROM results/ — DO NOT hand-edit.
[Placeholder: report the Kendall-$\tau$ between the token-matched and byte-matched rankings, and the
number of rank inversions, from Table~\ref{tab:rank} once the pilot sweep completes.]
Our results are complementary to concurrent work showing rankings are unstable under query
visibility~\citep{luo2026queryvisibility}: we show they are also unstable under the resource unit the
field measures them in. Code and all result files are released.
\end{abstract}

\section{Introduction}
\label{sec:intro}

% The problem: the field budgets in tokens, reports in memory. State the three verified facts as
% the motivation, then the contributions.
[Placeholder prose. Motivation: token-fraction budgeting is the de facto standard (KVPress
\texttt{compression\_ratio}; ratio/count budgets in \citealt{luo2026queryvisibility,
garcia2026protection, yuan2024longctxbench, li2025scbench}), yet the reported quantity is presented
as memory. Lead with the AdaKV masking finding (\S\ref{sec:audit}) as the existence proof.]

\paragraph{Contributions.}
\begin{enumerate}
  \item A three-way taxonomy of compression implementations---\emph{memory-faithful},
  \emph{simulation-only} (masked), and \emph{incommensurable} (off the token axis)---and a
  measurement instrument that assigns each method to a class by direct byte accounting
  (\S\ref{sec:method}).
  \item An audit (\S\ref{sec:audit}) showing the divergence between reported compression ratio and
  realized bytes across methods on the standard toolkit, including methods that free zero memory,
  and a first-principles account of why head-level allocation cannot help under GQA.
  \item A byte-matched re-ranking (\S\ref{sec:rerank}) of the standard method set, quantifying how
  the ordering changes when the budget is fixed in bytes rather than tokens, with a paired
  latency-at-matched-bytes measurement.
\end{enumerate}

% related_work.tex — \input into paper.tex, but written to read standalone.
% Organized by mechanism family; the argument builds toward Family 4 (adaptive allocation),
% which is the family a token-fraction budget cannot honestly express.

\section{Related Work}
\label{sec:related}

KV cache compression for autoregressive inference has grown into a large literature. We group it
by mechanism, and note for each family how it interacts with the resource unit --- token fraction
versus realized bytes --- that this paper argues is mis-specified throughout.

\paragraph{Static and heuristic eviction.}
The earliest training-free methods keep a fixed positional pattern. StreamingLLM~\cite{xiao2024streamingllm}
retains a few initial ``attention sink'' tokens plus a recent sliding window; LM-Infinite~\cite{han2024lminfinite}
independently proposes a $\Lambda$-shaped mask of the same form. These methods are memory-faithful
by construction --- what they keep is exactly what they store --- and, as recent audits find, they
are a far stronger baseline than the literature has credited~\cite{luo2026queryvisibility,garcia2026protection}.

\paragraph{Attention-score-based eviction.}
A second family scores tokens by accumulated attention and evicts the low-scoring ones. H2O~\cite{zhang2023h2o}
retains ``heavy hitters''; Scissorhands~\cite{liu2023scissorhands} exploits the persistence of
importance across steps; TOVA~\cite{oren2024tova} drops the token least attended by the most recent
query; FastGen~\cite{ge2024fastgen} profiles each head and applies a per-head policy; Keyformer~\cite{adnan2024keyformer}
de-biases the discarded-token distribution with a Gumbel-softmax score. A recurring critique is that
attention score is an incomplete importance signal --- value-vector norms matter too --- and that
these methods require the materialized attention matrix, making them incompatible with
FlashAttention~\cite{park2025keydiff}.

\paragraph{Query-aware / observation-window selection.}
SnapKV~\cite{li2024snapkv} introduced the dominant modern idea: pool attention from a short
\emph{observation window at the end of the prompt} to decide which prefix tokens to keep.
PyramidKV~\cite{cai2024pyramidkv} and PyramidInfer~\cite{zhang2024pyramidinfer} add a layer-wise
budget on top. The observation window is precisely the mechanism that concurrent audits identify as
confounded: because long-context benchmarks place the question last, the window overlaps the query,
and the method's apparent skill partly measures query--content overlap rather than intrinsic token
importance~\cite{luo2026queryvisibility}. Follow-up work treats the window as a defect to be fixed
(Judge~Q trains soft query tokens~\cite{liu2025judgeq}), reformulates selection to be query-agnostic
(KVzip~\cite{jang2025kvzip}), or removes the query from the score entirely (KeyDiff~\cite{park2025keydiff}).

\paragraph{Adaptive head- and layer-level budget allocation.}
This family is central to our argument. Rather than a uniform per-head budget, Ada-KV~\cite{feng2025adakv}
reallocates a global budget across heads by an $L_1$ eviction-loss bound; CAKE~\cite{qin2025cake}
allocates across layers by an attention-dispersion preference; HeadKV~\cite{fu2025headkv} scores
heads by retrieval-and-reasoning ability; RazorAttention~\cite{tang2025razorattention} and
DuoAttention~\cite{xiao2025duoattention} keep full cache only for identified retrieval/streaming
heads. Two facts make this family impossible to price in token fractions. First, under grouped-query
attention the unit of KV memory is the (layer, KV-group) cell, not the head: query heads in a group
share K and V, so \emph{head-level} budget concentration frees no memory unless the storage layout is
made ragged, which requires paged KV management~\cite{rehg2024kvcompress}. Second, some widely-used
implementations realize head-wise budgets by \emph{masking} rather than eviction --- the K/V tensors
are never shrunk --- so the reported compression ratio and the realized byte savings diverge
completely. We quantify both effects in \S\ref{sec:audit}.

\paragraph{Merging rather than eviction.}
CaM~\cite{zhang2024cam} merges to-be-evicted value caches into survivors; D2O~\cite{wan2025d2o}
adds cosine-similarity-weighted token merging under a per-layer dynamic budget; KVMerger~\cite{wang2024kvmerger}
clusters merge candidates by key similarity. Merging changes what a ``retained token'' means and
further loosens the token-count-to-bytes correspondence.

\paragraph{Quantization: a composable, orthogonal axis.}
KV quantization (KIVI~\cite{liu2024kivi}, KVQuant~\cite{hooper2024kvquant}) reduces bytes-per-element
rather than token count, and composes with eviction rather than competing with it. We do not evaluate
quantization, but it is the clearest evidence for our thesis: once precision varies, only a
\emph{byte} budget is well defined, and indeed the work nearest to ours in spirit trades token count
against bits-per-token and is forced to reason about memory rather than token fraction~\cite{zhang2024tokenprecision}.

\paragraph{Retrieval and offloading.}
Quest~\cite{tang2024quest}, InfLLM~\cite{xiao2024infllm}, and PQCache~\cite{zhang2025pqcache} never
permanently discard KV; they page or retrieve it, trading GPU bytes for latency and host memory.
These change the resource being budgeted and are out of our eviction-only scope, but they underline
that ``compression ratio'' means different things across families.

\paragraph{Evaluation and critique.}
The benchmark side is where this paper lives. longctx\_bench~\cite{yuan2024longctxbench} and
SCBench~\cite{li2025scbench} compare methods across families under compression \emph{ratios};
\citet{gao2025rethinking} and \citet{agrawal2026benchmarking} argue that ``compression ratio alone is
a poor predictor of end-to-end performance,'' and \citet{chen2026pitfalls} show that compression
damages some prompt content far more than aggregate scores reveal. Most directly,
\citet{luo2026queryvisibility} run a matched-\emph{token} audit and find that rankings invert under
query visibility, with SnapKV losing to a start-plus-window baseline. Our work is complementary and
concurrent: we hold their finding fixed and ask a different question --- whether the rankings are also
unstable under the \emph{resource unit}, correcting token fractions to realized bytes --- and we
localize a mechanism (masking-vs-eviction, GQA head sharing) inside the field's standard tooling.
Evaluation uses RULER~\cite{hsieh2024ruler}, LongBench~\cite{bai2024longbench}, and where possible
HELMET~\cite{yen2025helmet}.

\paragraph{The moving frontier.}
Finally, architectural methods increasingly make post-hoc eviction partly moot: native sparse
attention~\cite{yuan2025nsa}, mixture-of-block attention~\cite{lu2025moba}, and multi-head latent
attention~\cite{deepseek2024v2} learn or restructure sparsity rather than imposing it at inference.
These change \emph{what a KV byte is}, which is exactly why a byte-denominated evaluation protocol,
rather than a token-denominated one, is the durable way to compare across the field's future.


\section{Method: Byte Accounting for KV Compression}
\label{sec:method}

\paragraph{Realized bytes.}
For a cache with per-layer key/value tensors $K^{(\ell)}, V^{(\ell)}$ under method $m$,
\begin{equation}
  \mathrm{bytes}(m) \;=\; \underbrace{\sum_{\ell} \big(|K^{(\ell)}| + |V^{(\ell)}|\big)\,s_{\mathrm{dtype}}}_{\text{payload}}
  \;+\; \underbrace{\mathrm{meta}(m)}_{\text{accumulators, masks, tables}},
\label{eq:bytes}
\end{equation}
where $|\cdot|$ counts elements and $s_{\mathrm{dtype}}$ is bytes per element. Crucially, $|K^{(\ell)}|$
is measured from the \emph{stored tensor}, not inferred from a nominal token count: a method that masks
without slicing has full-size $K^{(\ell)}$ and thus full payload regardless of its reported ratio.
$\mathrm{meta}(m)$ is implemented per method and asserted non-zero wherever the method maintains state.

\paragraph{Sanity control.}
For full cache, Eq.~\eqref{eq:bytes} must reproduce the analytic
$2 \cdot n_{\text{layers}} \cdot n_{\text{kv-heads}} \cdot d_{\text{head}} \cdot n_{\text{tokens}} \cdot s_{\mathrm{dtype}}$
to within allocator granularity. \emph{Positive control:} \texttt{AdaKVPress} must measure at
$\approx$100\% of full cache at every nominal ratio; if it does not, the instrument---not the
hypothesis---is wrong.

\paragraph{Byte-matched budgeting.}
\texttt{ByteBudgetPress} wraps any method and binary-searches its nominal ratio, per
(method, model, context length), to hit a target realized-byte budget, so that all methods are
compared at equal bytes rather than equal token fraction.

\section{The Audit: Ratios Are Not Bytes}
\label{sec:audit}

\datafig{fig1_mispricing}{0.72}{%
  Nominal compression ratio (x) versus realized KV bytes (y), one line per method, one model, one
  context length. The dashed diagonal is what the ratio promises. Simulation-only methods appear as
  flat lines at full-cache bytes; sink/observation-window methods diverge upward by their reserved
  allocation. Generated from real measurements only.}{fig:mispricing}

\datafig{fig3_decomposition}{0.9}{%
  Per-method byte decomposition at a fixed nominal ratio: K/V payload, sink tokens, recency window,
  observation window, metadata, and peak-prefill transient.}{fig:decomp}

[Placeholder analysis. Report, from real runs: which methods land in each taxonomy class; the
realized-bytes value for AdaKV vs.\ its nominal ratio; the GQA (layer, KV-group) argument instantiated
with the evaluated model's head configuration.]

\section{Byte-Matched Re-Ranking}
\label{sec:rerank}

\datafig{fig2_reranking}{0.98}{%
  The correction. \textbf{Left:} accuracy versus nominal compression ratio (the field's view).
  \textbf{Right:} accuracy versus realized bytes (corrected). Rank inversions are annotated.}{fig:rerank}

\datafig{fig5_latency}{0.72}{%
  Accuracy per byte per second at matched bytes: methods that free no bytes also yield no decode
  speedup, while being ranked as if they saved both.}{fig:latency}

\begin{table}[t]
  \centering
  \caption{Kendall $\tau$ between the token-matched and byte-matched method rankings, per budget
  level (mean $\pm$ s.d. over seeds). $\tau < 1$ is the quantitative form of the re-ranking claim.
  Filled from \texttt{scripts/table1\_rank.py}; empty until the pilot completes.}
  \label{tab:rank}
  \begin{tabular}{lcccc}
    \toprule
    Budget (bytes) & B1 & B2 & B3 & B4 \\
    \midrule
    Kendall $\tau$ & --- & --- & --- & --- \\
    \# inversions  & --- & --- & --- & --- \\
    \bottomrule
  \end{tabular}
\end{table}

[Placeholder analysis of Table~\ref{tab:rank} and Figs.~\ref{fig:rerank}, \ref{fig:latency}.]

\section{Limitations}
\label{sec:limitations}
Single toolkit (NVIDIA KVPress) --- the finding is about how the field's standard implementations
account for memory, not a claim that every method's algorithm is flawed. Masking is a legitimate way
to study accuracy; the defect is the field's failure to distinguish simulation-faithful from
memory-faithful implementations. Pilot model/benchmark coverage is intentionally narrow; the extended
version broadens it. Allocator granularity puts a floor on byte-measurement precision, reported
explicitly. Concurrent work~\citep{luo2026queryvisibility} covers the query-visibility axis, which we
do not re-claim.

\section{Conclusion}
\label{sec:conclusion}
[Placeholder. Bytes, not ratios: correcting the resource unit changes the comparison, and in the
limit exposes methods that compress on paper but not in memory.]

\bibliographystyle{plainnat}
\bibliography{refs}

\end{document}
